\documentclass[11pt,a4paper]{article}

% ==============================================================================
% PREAMBLE: MODERN TECHNICAL ARTICLE STYLE (latex-notes-generator)
% ==============================================================================

% 1. Page Geometry & Typography
\usepackage[margin=2cm]{geometry}
\usepackage{mathpazo}
\usepackage{microtype}
\usepackage{setspace}
\setstretch{1.15}

% 2. Math & Symbols
\usepackage{amsmath, amssymb}

% 3. List spacing & Tables
\usepackage{enumitem}
\setlist{itemsep=0.3em, topsep=0.3em}
\usepackage{tabularx}
\usepackage{array}

% 4. Callout boxes (tcolorbox)
\usepackage[many]{tcolorbox}
\newtcolorbox{definitionbox}[1][]{
    colback=blue!3!white,
    colframe=blue!60!black,
    fonttitle=\bfseries,
    title=#1,
    arc=2mm
}
\newtcolorbox{formulabox}{
    colback=gray!5!white,
    colframe=gray!40!black,
    boxrule=0.5pt,
    arc=1mm
}

% 5. Header & Footer
\usepackage{fancyhdr}
\pagestyle{fancy}
\fancyhf{}
\fancyhead[L]{\textbf{Cumulative ML Code Glossary \& Practical Reference}}
\fancyhead[R]{ML \& Data Mining --- UniBo}
\fancyfoot[C]{\thepage}
\renewcommand{\headrulewidth}{0.4pt}

% 6. Hyperlinks
\usepackage{hyperref}
\hypersetup{
    colorlinks=true,
    linkcolor=blue!70!black,
    urlcolor=blue!70!black,
    citecolor=green!50!black
}

% 7. Code listings configuration
\usepackage{xcolor}
\usepackage{listings}
\usepackage{textcomp}

\definecolor{codegreen}{rgb}{0.12, 0.55, 0.15}
\definecolor{codegray}{rgb}{0.45, 0.45, 0.45}
\definecolor{codepurple}{rgb}{0.55, 0.12, 0.65}
\definecolor{codeblue}{rgb}{0.05, 0.25, 0.75}
\definecolor{codebase}{rgb}{0.55, 0.10, 0.25}      % Distinct berry/wine color for inline code & tokens
\definecolor{backcolour}{rgb}{0.96, 0.97, 0.99}    % Subtle, clean cool-gray background
\definecolor{boxframe}{rgb}{0.78, 0.84, 0.90}      % Slate frame border

\lstdefinestyle{moderncode}{
    backgroundcolor=\color{backcolour},
    commentstyle=\color{codegreen}\itshape,
    keywordstyle=\color{codeblue}\bfseries,
    numberstyle=\tiny\color{codegray},
    stringstyle=\color{codepurple},
    identifierstyle=\color{codebase},
    basicstyle=\ttfamily\scriptsize\color{codebase}, % Distinct berry/wine code color
    breaklines=true,
    breakatwhitespace=false,
    showstringspaces=false,
    keepspaces=true,
    numbers=left,
    numbersep=5pt,
    tabsize=4,
    frame=single,
    rulecolor=\color{boxframe},
    xleftmargin=1.2em,
    framexleftmargin=1.2em,
    aboveskip=0.8em,
    belowskip=0.6em
}
\lstset{style=moderncode}

% Distinct styling for \texttt to match code tokens
\let\oldtexttt\texttt
\renewcommand{\texttt}[1]{{\color{codebase}\oldtexttt{#1}}}

\begin{document}

% ==============================================================================
% GLOSSARY HEADER BANNER & TOC REGISTRATION (FUNCTIONAL HYPERREF LINK)
% ==============================================================================
\phantomsection
\addcontentsline{toc}{section}{\textbf{Cumulative ML Code Glossary \& Practical Reference}}

\begin{tcolorbox}[
    colback=blue!5!white,
    colframe=blue!65!black,
    arc=3mm,
    boxrule=1.2pt,
    left=14pt, right=14pt, top=12pt, bottom=12pt
]
    {\Huge\bfseries\color{blue!75!black} Cumulative ML Code Glossary}\\[0.3em]
    {\LARGE\bfseries Master Practical Playbook: What it Does, How \& When to Use It}\\[0.4em]
    {\normalsize\textsf{End-of-Handbook Cumulative Quick-Reference --- Expanding Laboratory by Laboratory}}\par
    \vspace{0.8em}
    \hrule
    \vspace{0.6em}
    \footnotesize
    \textbf{Course:} Machine Learning \& Data Mining (Prof. Claudio Sartori) \hfill \textbf{Degree:} Master in AI (UniBo)\\
    \textbf{Scope:} Centralized Method Reference, Practical Mechanics \& Production Pipeline Skeletons
\end{tcolorbox}
\vspace{0.5em}

This cumulative appendix provides a direct, practical playbook of all essential operations in Pandas, NumPy, Matplotlib, and Seaborn introduced across the laboratory sessions. For every technique, it breaks down:
\begin{itemize}[leftmargin=*]
    \item \textbf{Command / Syntax (How):} The clean, idiomatic Python call.
    \item \textbf{What it Does:} Exact technical action on the data structures in memory.
    \item \textbf{When to Use:} The analytical trigger and purpose in the machine learning workflow.
\end{itemize}

% ==============================================================================
% G.1 TABLE INGESTION, CARDINALITY & MEMORY DECOUPLING
% ==============================================================================
\phantomsection
\addcontentsline{toc}{subsection}{G.1 Table Ingestion, Cardinality \& Memory Decoupling}
\section*{G.1 Table Ingestion, Cardinality \& Memory Decoupling}

{\small
\noindent
\begin{tabularx}{\textwidth}{|>{\raggedright\arraybackslash\ttfamily\footnotesize\color{codebase}}p{5.8cm}|>{\raggedright\arraybackslash}X|>{\raggedright\arraybackslash}X|}
\hline
\textbf{\textcolor{black}{\textrm{Command / Syntax}}} & \textbf{What it Does} & \textbf{When to Use (Pipeline Context)} \\ \hline
pd.read\_csv(f, names=[...]) & Ingests raw CSV and assigns an explicit schema list. & When loading headerless files to avoid losing the first row as column titles. \\ \hline
df.copy() & Allocates an independent deep memory clone of the DataFrame. & Immediately after ingestion to protect raw baseline tables from in-place mutations. \\ \hline
df.shape / df.info() & Inspects dimensionality $(N, D)$, null counts, and storage dtypes. & At the first step of EDA to assess sample volume and detect float/object conversions. \\ \hline
df.nunique() & Counts distinct unique values per column. & At early screening to detect ID columns ($nunique=N$) and constants ($nunique \le 1$). \\ \hline
df.duplicated().sum() & Counts identical duplicate rows across all columns. & During data quality auditing to detect redundant records before model training. \\ \hline
assert condition, "msg" & Halts execution with an error if condition evaluates to False. & At the end of the pipeline to verify that data cleaning never mutated the raw table. \\ \hline
\end{tabularx}
}

\vspace{0.5em}
\noindent\textbf{Recurring Skeleton: Initial Data Hygiene \& Cardinality Audit}\par\vspace{0.2em}\nopagebreak
\begin{lstlisting}[language=Python]
# 1. Ingest with explicit schema and decouple memory
raw_df = pd.read_csv("dataset.csv")
df = raw_df.copy()

# 2. Inspect dimensions and audit cardinality
print(f"Dimensions: {df.shape[0]} rows, {df.shape[1]} columns")
cardinality = df.nunique()
candidate_ids = cardinality[cardinality == len(df)].index.tolist()
zero_variance = cardinality[cardinality <= 1].index.tolist()

# 3. Clean drop: remove IDs (memorization trap) and constants (zero entropy)
df = df.drop(columns=candidate_ids + zero_variance)
print(f"Dropped identifiers: {candidate_ids} | Dropped constants: {zero_variance}")
\end{lstlisting}

% ==============================================================================
% G.2 TYPE SANITIZATION & STEVENS' MEASUREMENT SCALES
% ==============================================================================
\phantomsection
\addcontentsline{toc}{subsection}{G.2 Type Sanitization \& Stevens' Measurement Scales}
\section*{G.2 Type Sanitization \& Stevens' Measurement Scales}

{\small
\noindent
\begin{tabularx}{\textwidth}{|>{\raggedright\arraybackslash\ttfamily\footnotesize\color{codebase}}p{5.8cm}|>{\raggedright\arraybackslash}X|>{\raggedright\arraybackslash}X|}
\hline
\textbf{\textcolor{black}{\textrm{Command / Syntax}}} & \textbf{What it Does} & \textbf{When to Use (Pipeline Context)} \\ \hline
s.astype("string") & Casts column to dedicated null-tolerant text storage. & When handling nominal codes (zip codes, store IDs) to disable meaningless numeric math. \\ \hline
s.str.replace(r"\textbackslash{}.0\$", "") & Strips trailing floating-point decimal string artifacts. & After integer columns containing nulls were coerced to float strings (\texttt{40100.0} $\to$ \texttt{"40100"}). \\ \hline
pd.Categorical(s, ordered=True) & Instantiates ordered category dtype backed by integer codes. & On ordinal features to enforce business ranking ($\text{Basic} < \text{Silver} < \text{Gold}$) over alphabetical order. \\ \hline
df.select\_dtypes("number") & Filters DataFrame to retain only numerical columns. & Before computing correlation heatmaps or distance metrics to exclude IDs and text codes. \\ \hline
\end{tabularx}
}

\vspace{0.5em}
\noindent\textbf{Recurring Skeleton: Enforcing Semantic Scales}\par\vspace{0.2em}\nopagebreak
\begin{lstlisting}[language=Python]
# 1. Fix float coercion on nominal codes (e.g. postal codes, terminal IDs)
df["postcode"] = df["postcode"].astype("string").str.replace(r"\.0$", "", regex=True)

# 2. Impose natural category order (avoids alphabetical sorting: Basic < Silver < Gold)
df["tier"] = pd.Categorical(
    df["tier"],
    categories=["Basic", "Silver", "Gold"],
    ordered=True
)
\end{lstlisting}

% ==============================================================================
% G.3 SENTINEL UNMASKING & MISSINGNESS DIAGNOSTICS (MCAR VS MNAR)
% ==============================================================================
\phantomsection
\addcontentsline{toc}{subsection}{G.3 Sentinel Unmasking \& Missingness Analysis (MCAR vs MNAR)}
\section*{G.3 Sentinel Unmasking \& Missingness Analysis (MCAR vs MNAR)}

{\small
\noindent
\begin{tabularx}{\textwidth}{|>{\raggedright\arraybackslash\ttfamily\footnotesize\color{codebase}}p{5.8cm}|>{\raggedright\arraybackslash}X|>{\raggedright\arraybackslash}X|}
\hline
\textbf{\textcolor{black}{\textrm{Command / Syntax}}} & \textbf{What it Does} & \textbf{When to Use (Pipeline Context)} \\ \hline
(s == sentinel).astype(int) & Creates a binary $0/1$ indicator feature from dummy codes. & \textit{Before} replacing sentinels with NaN, preserving the behavioral event (e.g. never bought). \\ \hline
s.replace(sentinel, np.nan) & Replaces dummy values with IEEE NaN out-of-place. & To clean numeric scales and unbias means and variances distorted by sentinel numbers. \\ \hline
s.isna().mean() * 100 & Calculates percentage of missing values in the column. & To evaluate data completeness and decide between imputation or column dropping. \\ \hline
df.groupby(s.isna())[y].mean() & Computes average target rate for observed vs missing rows. & To diagnose MCAR (rates match $\implies$ safe for median) vs MNAR (rates diverge $\implies$ keep flag). \\ \hline
\end{tabularx}
}

\vspace{0.5em}
\noindent\textbf{Recurring Skeleton: Sentinel Unmasking \& MNAR Diagnostic Pipeline}\par\vspace{0.2em}\nopagebreak
\begin{lstlisting}[language=Python]
# 1. Capture behavioral meaning of database sentinels (e.g. 999 = never purchased)
df["never_purchased"] = (df["recency_days"] == 999).astype(int)

# 2. Out-of-place replacement: unskew mean and standard deviation
df["recency_days"] = df["recency_days"].replace(999, np.nan)
df["age"] = df["age"].replace(0, np.nan)

# 3. Check if missingness correlates with the target (MNAR diagnostic)
for col in ["recency_days", "web_activity", "age"]:
    target_by_na = df.groupby(df[col].isna())["churn"].mean()
    divergence = abs(target_by_na.get(True, 0) - target_by_na.get(False, 0))
    # If divergence is substantial -> MNAR: keep binary indicator flag!
    print(f"Feature {col:15s} Missingness Divergence: {divergence:.3f}")
\end{lstlisting}

% ==============================================================================
% G.4 STATISTICAL ROBUSTNESS & CONTEXTUAL OUTLIER CONDITIONING
% ==============================================================================
\phantomsection
\addcontentsline{toc}{subsection}{G.4 Statistical Robustness \& Outlier Conditioning}
\section*{G.4 Statistical Robustness \& Outlier Conditioning}

{\small
\noindent
\begin{tabularx}{\textwidth}{|>{\raggedright\arraybackslash\ttfamily\footnotesize\color{codebase}}p{5.8cm}|>{\raggedright\arraybackslash}X|>{\raggedright\arraybackslash}X|}
\hline
\textbf{\textcolor{black}{\textrm{Command / Syntax}}} & \textbf{What it Does} & \textbf{When to Use (Pipeline Context)} \\ \hline
s.quantile([0.25, 0.75]) & Computes 25th ($Q_1$) and 75th ($Q_3$) percentiles. & To calculate robust dispersion metrics ($IQR = Q_3 - Q_1$) resistant to heavy tails. \\ \hline
np.log10(s.replace(0, nan)) & Applies point-wise base-10 logarithmic compression. & On right-skewed positive features (spend, counts) to stabilize variance for linear models. \\ \hline
df[s > upper\_fence] & Filters rows exceeding Tukey's fence ($Q_3 + 1.5 \cdot IQR$). & To isolate candidate outliers for inspection before deciding whether to drop or keep. \\ \hline
df.boxplot(col, by=group) & Draws stratified boxplots across category levels. & To check whether extreme values belong to legitimate VIP segments (e.g., Gold whales). \\ \hline
\end{tabularx}
}

\vspace{0.5em}
\noindent\textbf{Recurring Skeleton: Tukey IQR Boundaries \& Subgroup Inspection}\par\vspace{0.2em}\nopagebreak
\begin{lstlisting}[language=Python]
# 1. Establish Tukey boundaries
q1 = df["revenue"].quantile(0.25)
q3 = df["revenue"].quantile(0.75)
iqr = q3 - q1
upper_bound = q3 + 1.5 * iqr

# 2. Isolate candidate outliers
candidates = df[df["revenue"] > upper_bound]

# 3. Check distribution across business segments before deleting!
print("Extreme values breakdown across tiers:")
print(candidates["tier"].value_counts(normalize=True))
# Rule: If concentrated in high-value segment (e.g. Gold whales), do not delete!
\end{lstlisting}

% ==============================================================================
% G.5 TARGET IMBALANCE & MAJORITY BASELINE COMPARISON
% ==============================================================================
\phantomsection
\addcontentsline{toc}{subsection}{G.5 Target Imbalance \& Majority Baseline Comparison}
\section*{G.5 Target Imbalance \& Majority Baseline Comparison}

{\small
\noindent
\begin{tabularx}{\textwidth}{|>{\raggedright\arraybackslash\ttfamily\footnotesize\color{codebase}}p{5.8cm}|>{\raggedright\arraybackslash}X|>{\raggedright\arraybackslash}X|}
\hline
\textbf{\textcolor{black}{\textrm{Command / Syntax}}} & \textbf{What it Does} & \textbf{When to Use (Pipeline Context)} \\ \hline
s.value\_counts(normalize=True) & Computes relative class proportions summing to 1.0. & On the target column to quantify imbalance severity ($88/12$) before modeling. \\ \hline
pd.Series(val, index=...) & Creates a constant Series aligned with dataset index. & To construct trivial baseline predictions ($y_{\text{pred}} = \text{mode}(y)$) for benchmarking. \\ \hline
pd.crosstab(actual, pred) & Computes a 2D contingency table of actual vs predicted classes. & To unmask false performance: shows 0 true positives despite high apparent accuracy. \\ \hline
((y == 1) \& (pred == 1)).sum() & Counts True Positives via bitwise AND reduction. & To evaluate whether a model actually catches the rare class of interest (Recall). \\ \hline
ax.axhline(overall\_rate) & Draws a horizontal scalar reference line across plot axes. & On categorical rate bar charts to show whether subgroups deviate from the baseline average. \\ \hline
\end{tabularx}
}

\vspace{0.5em}
\noindent\textbf{Recurring Skeleton: The Majority-Class Baseline Check}\par\vspace{0.2em}\nopagebreak
\begin{lstlisting}[language=Python]
actual = df["churned"]
majority_class = actual.mode()[0]
trivial_pred = pd.Series(majority_class, index=df.index)

# 2x2 contingency table showing deceptive accuracy
conf_mat = pd.crosstab(actual, trivial_pred, rownames=["Actual"], colnames=["Pred"])
acc = (actual == trivial_pred).mean()
actual_positives = (actual == 1).sum()
true_positives = ((actual == 1) & (trivial_pred == 1)).sum()
recall = true_positives / actual_positives if actual_positives > 0 else 0.0

print("Contingency Matrix:\n", conf_mat)
print(f"Majority Baseline Accuracy: {acc:.1%} | Minority Recall: {recall:.1%}")
# Accuracy is misleading: on imbalanced data, evaluate Recall and Precision!
\end{lstlisting}

% ==============================================================================
% G.6 BIVARIATE ASSOCIATIONS, HEATMAPS & SIMPSON'S PARADOX
% ==============================================================================
\phantomsection
\addcontentsline{toc}{subsection}{G.6 Bivariate Associations, Heatmaps \& Simpson's Paradox}
\section*{G.6 Bivariate Associations, Heatmaps \& Simpson's Paradox}

{\small
\noindent
\begin{tabularx}{\textwidth}{|>{\raggedright\arraybackslash\ttfamily\footnotesize\color{codebase}}p{5.8cm}|>{\raggedright\arraybackslash}X|>{\raggedright\arraybackslash}X|}
\hline
\textbf{\textcolor{black}{\textrm{Command / Syntax}}} & \textbf{What it Does} & \textbf{When to Use (Pipeline Context)} \\ \hline
df[cols].corr("pearson") & Computes linear correlation matrix $r \in [-1, 1]$. & To detect linear associations and redundant pairs ($|r| > 0.85$) on numeric features. \\ \hline
df[cols].corr("spearman") & Computes rank-order correlation matrix $\rho \in [-1, 1]$. & When relationships may be non-linear or features contain heavy tails. \\ \hline
sns.heatmap(corr, annot=True) & Draws a color-coded matrix with numeric annotations. & To quickly audit collinearity across all numeric predictors at a glance. \\ \hline
for k, g in df.groupby(by): & Iterates over stratified subgroup DataFrame slices. & To detect Simpson's Paradox: tests whether aggregate correlations hold inside subgroups. \\ \hline
\end{tabularx}
}

\vspace{0.5em}
\noindent\textbf{Recurring Skeleton: Simpson's Paradox Check (Grouped Correlations)}\par\vspace{0.2em}\nopagebreak
\begin{lstlisting}[language=Python]
# 1. Aggregate correlation across the entire dataset
overall_r = df["spend"].corr(df["web_sessions"])
print(f"Aggregate Correlation: {overall_r:.2f}")

# 2. Stratified correlation within each categorical level
print("Correlation within each subgroup:")
for level, group in df.groupby("tier", observed=False):
    intra_r = group["spend"].corr(group["web_sessions"])
    print(f"  Level {str(level):10s}: {intra_r:.2f}")

# Simpson's Paradox detected if subgroup correlation collapses or reverses!
\end{lstlisting}

\end{document}
